\section*{Chapitre \ref{ch:ll:the-strategy-engine} --- Le moteur de stratégie}

\begin{solution}{exo:ll:the-strategy-engine:1}
$34\,200.012345~\text{s}$ correspond à $342\,000\,123$ intervalles de \qty{100}{\micro\second}, et $342\,000\,123 \bmod 1\,024 = 507$. La
roue couvre $1\,024 \times \qty{100}{\micro\second} = \qty{102.4}{\milli\second}$, donc un temporisateur situé à 5 secondes est à environ 48.8
tours d'écart~: il attend dans son emplacement pendant que la roue passe devant lui 48 fois.
\end{solution}

\begin{solution}{exo:ll:the-strategy-engine:2}
Par instant, puis identifiant~: (120, 7), (120, 9), (121, 3).
\end{solution}

\begin{solution}{exo:ll:the-strategy-engine:3}
En rejeu, l'instant de l'événement est celui où le moteur en production l'a vu, si bien que les décisions dépendent des mêmes nombres. Cela ne concorde avec l'exécution en production que
si le moteur en production tire lui aussi son temps des événements (ou d'une horloge dont les lectures sont enregistrées avec les entrées), et si
les entrées enregistrées sont exactement celles que le moteur en production a reçues, dans le même ordre.
\end{solution}

\begin{solution}{exo:ll:the-strategy-engine:4}
$k\,(1 - (1 - 1/k)^{100})$ en moyenne. Elle vaut 93 pour $k \approx 700$~: le résultat dépend de plusieurs centaines
d'ordres possibles (moins de cent ordres plafonneraient le compte à $k$~; des probabilités inégales en demandent davantage).
\end{solution}

\begin{solution}{exo:ll:the-strategy-engine:5}
La première sous les anciens paramètres (deux ticks), la seconde sous les nouveaux (quatre)~: le snapshot est installé avant le premier
événement survenant à son instant d'effet ou après. Le journal indique la version de chaque action~: la version 1 avant, la version 2 après.
\end{solution}

\begin{solution}{exo:ll:the-strategy-engine:6}
Oui~: $2 \times 10^6 \times \qty{53}{\nano\second} \approx 0.11$, soit un taux d'utilisation d'environ 11~\%, donc aucune file d'attente ne se forme (la queue
de \qty{130}{\nano\second} au p99 ne change rien à cela).
\end{solution}

\begin{solution}{exo:ll:the-strategy-engine:7}
Programmer un temporisateur de rafraîchissement à chaque cotation, à échéance dans \qty{5}{\second}, avec l'identifiant de la cotation dans celui du temporisateur~; quand il se déclenche et
la cotation est inchangée, émettre un replace au même prix et programmer le rafraîchissement suivant. Les deux moteurs doivent utiliser les mêmes
identifiants et la même règle de tri pour que les empreintes de rejeu concordent, et l'empreinte validée est régénérée une fois à partir du
moteur C++.
\end{solution}

\begin{solution}{exo:ll:the-strategy-engine:8}
L'horodatage diffère à chaque exécution, si bien que la sortie n'est jamais reproductible bit à bit, et tout code qui la compare à quoi que ce soit
(un délai d'expiration, un âge) fait aussi diverger les ordres~; l'horloge système peut aussi sauter lorsqu'on la corrige. Prendre l'horodatage de la piste d'audit
à partir de l'instant de l'événement dans le moteur, ou l'ajouter hors du moteur là où les messages sortent de la machine (chapitre~23), et le garder
hors de la comparaison rejouée.
\end{solution}

\begin{solution}{pb:ll:the-strategy-engine:1}
\begin{enumerate}
\item Une lecture de l'horloge murale~: 100 sorties distinctes~; un ordre de hachage à graine fixée~: 93~; deux threads sur une même file~: 49.
\item Chaque exécution lit des instants différents, et ces instants figurent dans la sortie.
\item Les sorties dépendent de l'ordre que produit l'ensemble, et plusieurs graines conduisent au même ordre~: plusieurs centaines d'ordres,
répétés de temps à autre sur cent exécutions.
\item Beaucoup d'entrelacements aboutissent à la même sortie lorsque les threads alternent par hasard comme ils l'ont fait auparavant~; davantage de contention ou
des consommateurs plus lents en donneraient davantage.
\item Journaliser la sortie des deux exécutions avec leurs numéros de séquence et les comparer enregistrement par enregistrement (ou procéder par dichotomie sur l'empreinte des préfixes).
\item L'horloge murale~: décisions identiques avec des horodatages différents~; l'ordre de hachage~: mêmes actions dans un ordre différent ou
avec des identifiants différents~; les threads~: un état de carnet que l'entrée n'a jamais produit.
\item La graine est tirée au démarrage, une fois par processus~: les relances dans un même processus la partagent.
\item L'entrelacement des threads, que le débogueur modifie en arrêtant les threads.
\item Tirer le temps des événements grâce à une \omterm{def:ll:the-strategy-engine:wheel}{horloge injectée}, et horodater hors du moteur.
\item Réserver les tables de hachage aux seules recherches, et itérer dans un ordre défini (un tableau des clés, une table ordonnée ou un tri) lorsque
l'ordre décide de quoi que ce soit.
\item Fusionner les deux entrées selon une règle (instant, puis source, puis séquence) dans un seul thread en amont du moteur, ou donner à chaque entrée son
propre anneau et laisser le moteur fusionner.
\item Un~; la firme en donnait un.
\item \textbf{Résultat nommé.} Sur cent relances des mêmes 6\,000 événements~: 100 résultats distincts avec une lecture de l'horloge murale, 93 avec
un ordre de hachage à graine fixée, 49 avec deux threads sur une même file, 100 avec les trois, et un après les corrections (et un pour le moteur
de la firme).
\item Un backtest non reproductible ne peut être ni débogué ni comparé~: une petite différence peut masquer un gros bug, et deux
versions d'une stratégie ne peuvent être distinguées du bruit.
\item Des réductions en virgule flottante dans un ordre différent, des options de compilation ou des bibliothèques différentes (chapitre~8), des données d'entrée différentes
(un paquet manqué dans un enregistrement), et une mémoire non initialisée.
\item Les opérations entières sont exactes et définies de la même manière dans les deux langages, alors que les résultats en virgule flottante dépendent de l'ordre
et de la contraction des opérations.
\item Rejouer une journée enregistrée deux fois et comparer les empreintes~; la rejouer sur chaque construction du moteur et la comparer à l'empreinte
validée~; exécuter le rejeu sous un sanitiseur.
\item Tout incident peut être rejoué à l'identique sur un desk, avec un débogueur, autant de fois que nécessaire.
\item Hors du chemin de décision~: supervision, journalisation des temps, outils opérationnels dont les sorties ne font pas partie de l'état
rejoué.
\item Un seul thread, des entrées ordonnées, un temps tiré des événements, et aucune itération selon un ordre qui varie.
\end{enumerate}
\end{solution}

\begin{solution}{iq:ll:the-strategy-engine:1}
Aucun verrou, aucune coordination, l'état dans les caches d'un seul cœur, et un ordre unique des événements qui rend le moteur déterministe
et rejouable. Le coût~: le débit d'un seul cœur constitue la limite, un gestionnaire lent retarde tous les autres événements, et le travail qui pourrait s'exécuter
en parallèle (analyse, journalisation) doit être déplacé vers d'autres threads par des anneaux.
\iqlookfor{déterminisme et localité face à la capacité d'un seul cœur.}
\end{solution}

\begin{solution}{iq:ll:the-strategy-engine:2}
Une \omterm{def:ll:the-strategy-engine:wheel}{roue de temporisateurs} hachée~: programmation et annulation en temps constant, emplacements parcourus à mesure que le temps avance, une roue hiérarchique ou
des compteurs de tours pour les longs temporisateurs~; déclencher les temporisateurs échus dans un ordre défini (instant, puis identifiant).
\iqlookfor{la roue, et l'ordre pour le déterminisme.}
\end{solution}

\begin{solution}{iq:ll:the-strategy-engine:3}
D'abord les entrées (les événements enregistrés sont-ils ceux que la production a reçus, dans le même ordre~?), puis le temps (du code lit-il l'horloge
au lieu de l'instant de l'événement~?), puis l'ordre (threads, itération sur les tables de hachage), puis les aspects numériques (ordre des opérations en virgule flottante, options de compilation), puis
les paramètres (les versions en vigueur). Comparer les deux sorties enregistrement par enregistrement pour trouver la première différence.
\iqlookfor{entrées, temps, ordre, numérique, paramètres.}
\end{solution}

\begin{solution}{iq:ll:the-strategy-engine:4}
Publier des snapshots complets et versionnés sur un anneau de contrôle, avec un instant d'effet~; le moteur en installe un entre deux événements, jamais
pendant un événement, et horodate chaque action avec la version en vigueur~; en rejeu, les snapshots sont des entrées.
\iqlookfor{snapshots complets entre deux événements, et versionnement.}
\end{solution}

\begin{solution}{iq:ll:the-strategy-engine:5}
Allocations, appels système, verrous, attentes, mise en forme de texte, boucles non bornées et tout ce qui varie d'une exécution à l'autre. Vérifier à l'aide de
compteurs d'allocations dans les tests, traçage des appels système, empreintes de rejeu et \omterm{def:ll:where-latency-comes-from:tail}{latences de queue} mesurées par événement.
\iqlookfor{une liste et une manière de vérifier chaque point.}
\end{solution}

\begin{solution}{iq:ll:the-strategy-engine:6}
Un thread par moteur, alimenté par des anneaux~; entrées fusionnées selon une règle et enregistrées~; temps injecté (simulé en rejeu, réel en production)~; temporisateurs
sur une roue à déclenchement ordonné~; paramètres sous forme de snapshots versionnés entre deux événements~; aucune dépendance à l'ordre de hachage ni à la virgule flottante
dans les décisions~; garanti par des tests de rejeu avec empreintes validées, en intégration continue et après chaque livraison.
\iqlookfor{règles et tests qui les font respecter.}
\end{solution}
